\section{Bottleneck Analysis}

\subsection{What still delays the first PV command}
\label{sec:current-headroom}

After QK completes, the first K64 PV half still depends on two N32 score
fragments. Each fragment must be loaded from TMEM, reduced, assigned an MX
scale, mapped and packed to E2M1, included in the represented denominator,
written into the score overlay, and published. Q0 alone is not useful; Q1
must reach the same point before the matrix issuer has a legal operand. The
other query stage and independent arithmetic hide part, but not all, of this
chain.

At B1/S4096/H24/D128, a clean tuning record measured 0.092448 ms and repeated
at 0.092512 ms. Three intentionally incorrect kernels bound the local
headroom:
\begin{table}[htbp]
\centering
\caption{Speed-of-light diagnostics for the retained GB200 schedule. These
rows do not compute valid attention.}
\label{tab:speed-of-light}
\small
\begin{tabular}{lrr}
\toprule
Diagnostic & Time (ms) & Gap from record \\
\midrule
score-pack ceiling & 0.091168 & 1.280\us\ (1.38\%) \\
row max retained, raw score pack & 0.090112 & 2.336\us\ (2.53\%) \\
fixed P & 0.087616 & 4.832\us\ (5.23\%) \\
\bottomrule
\end{tabular}
\end{table}
The fixed-P experiment removes nearly all P construction yet saves only
5.23\%. Early full-FP4 kernels were dominated by P quantization; the retained
NV/MX code path has removed most of that exposed cost. The remaining distance
from four-times BF16 also contains TMEM score loads, K64 publication
granularity, shared-memory delivery, two FP32 output accumulators, online
correction, and the epilogue. Another local polynomial fit cannot remove those
terms.

\subsection{Instruction count is not the schedule}

A matched fixed-P diagnostic cut dynamic instructions from about 54.6 million
to 11.9 million and raised tensor-pipe activity from 18.8\% to 26.6\%. Other
experiments removed arithmetic but became slower. Independent instructions
had occupied scheduler slots while a TMEM dependency or barrier event was
pending; removing them exposed the wait without advancing publication.

The represented-denominator decoder is a small example:
\begin{table}[htbp]
\centering
\caption{Equivalent denominator reductions under the same schedule.}
\label{tab:denom-schedule}
\small
\begin{tabular}{lrr}
\toprule
Decoder & Time (ms) & Change \\
\midrule
aggregate, retained & 0.092512 & baseline \\
pairwise partials & 0.094208 & +1.83\% \\
stream after each word & 0.094464 & +2.11\% \\
\bottomrule
\end{tabular}
\end{table}
All three are mathematically equivalent. The alternatives extend live ranges
and alter compiler interleaving. This also explains why an integer threshold
classifier lost to packed FMA plus native conversion despite looking simpler
algebraically.

\subsection{Why B300 is not uniformly 1.5 times faster}

\begin{table}[htbp]
\centering
\caption{Hardware properties relevant to the measured kernels. Per-SM
exponential rates are reported by FA4; GPU-level Ultra ratios are from
NVIDIA.}
\label{tab:blackwell-generations}
\small
\begin{tabular}{lcc}
\toprule
Property & GB200 / SM100 & B300 / SM103 \\
\midrule
Visible SMs in this study & 152 & 148 \\
Maximum reported clock & 2062 MHz & 2032 MHz \\
TMEM per SM & 256 KB & 256 KB \\
Key exponential rate & 16 ops/clock/SM & 32 ops/clock/SM \\
Dense NVFP4 GPU class & 1.0$\times$ & 1.5$\times$ \\
Fused TMEM load + reduction & no & \code{tcgen05.ld.red} \\
\bottomrule
\end{tabular}
\end{table}

Blackwell Ultra raises dense NVFP4 throughput by 1.5$\times$, doubles key SFU
rates, and exposes fused TMEM load/reduction
\cite{nvidia2025blackwellultra,nvidia2026ptx,zadouri2026flashattention4}.
Those changes do not scale the whole attention loop. Every CTA still occupies
all 512 TMEM columns, so residency remains one CTA per SM. Faster QK/PV makes
score reduction, scale publication, online state, and the epilogue a larger
fraction of time.

Launch geometry caused the first apparent regression. A binary compiled for
152 persistent workers ran on a 148-SM B300 and left four CTAs in an almost
empty second wave. Recompiling for 148 removes that tail. The retained D128
SM103 policy then uses six K/V stages, fused score reduction, and only a light
quarter-3 EX2 route where it wins. Routing EX2 through every quarter ties or
loses despite the stronger SFU. The full density and grid sweep is moved to
Appendix~\ref{app:sm103-tuning}.

At S8192/H64 and S9472/H64, enough work is available to sustain 3116 and 3159
TFLOP/s. S9472 has 2368 logical query-pair jobs, exactly 16 per persistent
CTA. At S32768/H24, the tested B300 instead reaches 2945 TFLOP/s versus 2998
on GB200. Its visible SM-clock envelope is
\BThreeSMClockEnvelopeDeficitPct\% smaller; normalized per SM-clock, B300 is
\BThreePerSMClockGainPct\% faster. The normalization explains the direction
but is not an observed throughput result.

\subsection{D64 specialization}
\label{sec:b300-d64}

HAO's published D64 full-FP4 route hits an atom boundary: its scaled FP4 PV
atom requires K128. Our D64 path uses dedicated K64 NVFP4/MXFP4 atoms. On
SM103 it also uses \code{tcgen05.ld.red} for all four score quarters and sends
eight of each quarter's 16 score pairs through native EX2.

Across all 24 matched shapes, safe NV/MX on B300 is
\BThreeDsixtyfourSafeGenerationGeo$\times$ faster than GB200 geometrically,
and \BThreeDsixtyfourSafeSaturatedGenerationGeo$\times$ for $S\ge4096$. At
S32768/H24 it reaches \BThreeDsixtyfourNVMXTflops\ TFLOP/s in
\BThreeDsixtyfourNVMXTime\,ms. A raw NV/NV control is faster on some inputs
but produces non-finite values at H32/S32768 for one seed; the bounded control
and full matrix are in Appendix~\ref{app:nvnv-d64}.

\subsection{What would move the boundary}

The remaining useful changes require a different hardware or ownership
contract:
\begin{description}
  \item[Scaled-FP4 K32 PV.] Consume each N32 completion immediately instead
  of waiting for a pair.
  \item[Another score bank.] Give QK a destination while PV still owns the
  current P overlay.
  \item[Scales outside TMEM.] Let the matrix instruction consume compact
  scales directly from shared memory or registers.
  \item[Compressed score ownership.] Create look-ahead without duplicating a
  full 128-column FP32 score tile.
  \item[Wider SM103 tiles.] K96/K64 Ultra instructions make N192 or N256
  attractive, but only if two-query K/V reuse and the P lifecycle survive the
  rewrite.
\end{description}
